% Part: first-order-logic
% Chapter: completeness
% Section: completeness-theorem

\documentclass[../../../include/open-logic-section]{subfiles}

\begin{document}

\iftag{FOL}
      {\olfileid{fol}{com}{cth}}
      {\olfileid{pl}{com}{cth}}

\olsection{సంపూర్ణతా సిద్ధాంతం}

\begin{explain}
మన ఫలితాలను కలుపుదాం: సంపూర్ణతా సిద్ధాంతానికి చేరుకుంటాం.
\end{explain}

\begin{thm}[సంపూర్ణతా సిద్ధాంతం]
\ollabel{thm:completeness}
$\Gamma$ ఒక \tetoken{వాక్యాలు}{!!{sentence}s} సమితి అనుకుందాం. $\Gamma$ అవైరుధ్యమైతే,
అది సంతృప్తిపరచదగినది.
\end{thm}

\begin{proof}
$\Gamma$ అవైరుధ్యమని అనుకుందాం.\iftag{FOL}{
  \olref[hen]{lem:henkin} ప్రకారం, $\Gamma' \supseteq \Gamma$ అయ్యే
  సంతృప్తీకృత అవైరుధ్య సమితి ఉంది.}{}
\olref[lin]{lem:lindenbaum} ప్రకారం,
$\Gamma^* \supseteq \iftag{FOL}{\Gamma'}{\Gamma}$ అయ్యే అవైరుధ్య,
\tetoken{సంపూర్ణ}{!!{complete}} సమితి ఉంది.
\iftag{FOL}{
  $\Gamma' \subseteq \Gamma^*$ కనుక, ప్రతి \tetoken{సూత్రం}{!!{formula}}~$!A(x)$కు
  $\Gamma^*$లో
  \iftag{prvEx}
        {$\lexists[x][!A(x)] \lif !A(c)$}
        {$\lnot\lforall[x][!A(x)] \lif \lnot !A(c)$}
  రూపంలోని \tetoken{ఒక వాక్యం}{!!a{sentence}} ఉంటుంది; అందువల్ల $\Gamma^*$
  సంతృప్తీకృతమైనది. $\Gamma$లో~$\eq$ లేకపోతే,}{అప్పుడు}
\olref[mod]{lem:truth} ప్రకారం,
$\iftag{FOL}{\Sat{M(\Gamma^*)}}{\pSat{v(\Gamma^*)}}{!A}$ అయినప్పుడు,
అప్పుడు మాత్రమే $!A \in \Gamma^*$. ప్రత్యేకంగా, ప్రతి $!A \in
\Gamma$కు $\iftag{FOL}{\Sat{M(\Gamma^*)}}{\pSat{v(\Gamma^*)}}{!A}$;
కాబట్టి $\Gamma$ సంతృప్తిపరచదగినది.\iftag{FOL}{
  $\Gamma$లో~$\eq$ ఉంటే, \olref[ide]{lem:truth} ప్రకారం అన్ని
  \tetoken{వాక్యాలు}{!!{sentence}s}~$!A$కు $\Sat{\equivclass{M}{\approx}}{!A}$ అయినప్పుడు,
  అప్పుడు మాత్రమే $!A \in \Gamma^*$. ప్రత్యేకంగా, ప్రతి $!A \in
  \Gamma$కు $\Sat{\equivclass{M}{\approx}}{!A}$; కాబట్టి $\Gamma$
  సంతృప్తిపరచదగినది.}{}
\end{proof}

\begin{cor}[సంపూర్ణతా సిద్ధాంతం, రెండవ రూపం]
\ollabel{cor:completeness}
అన్ని~$\Gamma$, \tetoken{వాక్యాలు}{!!{sentence}s}~$!A$లకు: $\Gamma \Entails !A$ అయితే
$\Gamma \Proves !A$.
\end{cor}

\begin{proof}
\olref{cor:completeness}, \olref{thm:completeness}లలోని~$\Gamma$లు
సార్వత్రికంగా పరిమాణీకృతమయ్యాయని గమనించండి. మనం గందరగోళపడకుండా,
\olref{thm:completeness}ను వేరొక చరరాశితో మళ్లీ చెబుదాం: ఏ
\tetoken{వాక్యాలు}{!!{sentence}s} సమితి~$\Delta$కైనా, $\Delta$ అవైరుధ్యమైతే అది
సంతృప్తిపరచదగినది. ప్రతిపక్షం ద్వారా, $\Delta$
సంతృప్తిపరచదగినది కాకపోతే $\Delta$ వైరుధ్యభరితమైనది. ఉపఫలితాన్ని
నిరూపించడానికి దీన్ని ఉపయోగిస్తాం.

$\Gamma \Entails !A$ అని అనుకుందాం. అప్పుడు
\olref[syn][sem]{prop:entails-unsat} ప్రకారం $\Gamma \cup \{\lnot !A\}$
సంతృప్తిపరచలేనిది. మన~$\Delta$గా $\Gamma \cup \{\lnot !A\}$ను
తీసుకుంటే, \olref{thm:completeness} యొక్క ముందరి రూపం $\Gamma \cup
\{\lnot !A\}$ వైరుధ్యభరితమని ఇస్తుంది. కింది ఫలితం ప్రకారం
\iftag{FOL}{%
  \tagrefs{prfAX/{fol:axd:prv:prop:prov-incons},
    prfSC/{fol:seq:prv:prop:prov-incons},
    prfND/{fol:ntd:prv:prop:prov-incons},
    prfTab/{fol:tab:prv:prop:prov-incons}}}{%
  \tagrefs{prfAX/{pl:axd:prv:prop:prov-incons},
    prfSC/{pl:seq:prv:prop:prov-incons},
    prfND/{pl:ntd:prv:prop:prov-incons},
    prfTab/{pl:tab:prv:prop:prov-incons}}},
$\Gamma \Proves !A$.
\end{proof}

\tagprob{FOL}
\begin{prob}
\olref[fol][com][cth]{cor:completeness}ను ఉపయోగించి
\olref[fol][com][cth]{thm:completeness}ను నిరూపించండి; తద్వారా
సంపూర్ణతా సిద్ధాంతపు రెండు సూత్రీకరణలు సమానార్థకాలని చూపండి.
\end{prob}
\tagendprob

\tagprob{notFOL}
\begin{prob}
\olref[pl][com][cth]{cor:completeness}ను ఉపయోగించి
\olref[pl][com][cth]{thm:completeness}ను నిరూపించండి; తద్వారా
సంపూర్ణతా సిద్ధాంతపు రెండు సూత్రీకరణలు సమానార్థకాలని చూపండి.
\end{prob}
\tagendprob

\tagprob{FOL}
\begin{prob}
ఒక \tetoken{వ్యుత్పత్తి}{!!a{derivation}} వ్యవస్థ సంపూర్ణంగా ఉండాలంటే, ప్రతి
సంతృప్తిపరచలేని సమితీ వైరుధ్యభరితమని నిరూపించగలిగేంత బలమైన నియమాలు
దానికి ఉండాలి. సంపూర్ణతను నిరూపించడానికి \tetoken{వ్యుత్పత్తి}{!!{derivation}} నియమాల్లో ఏవి
అవసరమయ్యాయి? ఈ నియమాల్లో ఏవైనా నిరూపణలో ఎక్కడా ఉపయోగించబడలేదా? ఈ
ప్రశ్నలకు సమాధానం ఇవ్వడానికి,
\olref[fol][com][cth]{thm:completeness} నిరూపణకు దారితీసే ఏ ఫలితాల్లో
ఏ \tetoken{వ్యుత్పత్తి}{!!{derivation}} నియమాలు ఉపయోగించబడ్డాయో చూపే జాబితా లేదా రేఖాచిత్రం
తయారుచేయండి. ఈ నిరూపణల్లో నియమాల అవ్యక్త ఉపయోగాలనూ తప్పకుండా
గుర్తించండి.
\end{prob}
\tagendprob

\tagprob{notFOL}
\begin{prob}
ఒక \tetoken{వ్యుత్పత్తి}{!!a{derivation}} వ్యవస్థ సంపూర్ణంగా ఉండాలంటే, ప్రతి
సంతృప్తిపరచలేని సమితీ వైరుధ్యభరితమని నిరూపించగలిగేంత బలమైన నియమాలు
దానికి ఉండాలి. సంపూర్ణతను నిరూపించడానికి \tetoken{వ్యుత్పత్తి}{!!{derivation}} నియమాల్లో ఏవి
అవసరమయ్యాయి? ఈ నియమాల్లో ఏవైనా నిరూపణలో ఎక్కడా ఉపయోగించబడలేదా? ఈ
ప్రశ్నలకు సమాధానం ఇవ్వడానికి,
\olref[pl][com][cth]{thm:completeness} నిరూపణకు దారితీసే ఏ ఫలితాల్లో
ఏ \tetoken{వ్యుత్పత్తి}{!!{derivation}} నియమాలు ఉపయోగించబడ్డాయో చూపే జాబితా లేదా రేఖాచిత్రం
తయారుచేయండి. ఈ నిరూపణల్లో నియమాల అవ్యక్త ఉపయోగాలనూ తప్పకుండా
గుర్తించండి.
\end{prob}
\tagendprob

\end{document}
